% A clean paper with what hyperref writes for links, bookmarks, and the start view, and an attached data file.
% Shared preamble of the active-content fixtures. The pdfTeX settings keep builds reproducible and keep the build
% path out of the PDF.
\pdftrailerid{}
\pdfinfoomitdate=1
\pdfsuppressptexinfo=-1
\documentclass[letterpaper,11pt]{article}
\usepackage[margin=1in]{geometry}
\pagestyle{empty}
\newcommand{\paperintro}{%
Flaky tests pass and fail on the same code without any change to the program under test.
This study reruns each test of 26 open-source Java projects in shuffled orders and records
every outcome together with the order in which the tests ran.}
\newcommand{\paperresults}{%
The detector labels 450 of 533 tests as stable, 51 as flaky, and 32 as unclear. A manual check
of the flaky tests confirms 47 of the 51 labels.}

\usepackage{embedfile}
\usepackage{hyperref}
\embedfile{tests.csv}
\begin{document}
\section{Introduction}\label{sec:intro}
\paperintro{} The data and the scripts are available at
\url{https://anonymous.4open.science/r/flaky-1234}, as Section~\ref{sec:results} and
prior work~\cite{luo} describe.

\section{Results}\label{sec:results}
\paperresults{} The definition of flakiness follows \href{https://example.org/flaky-tests}{prior work}.

\begin{thebibliography}{1}
\bibitem{luo} A. Author and B. Author, ``A study of flaky tests,'' in \emph{Proc. Conf.}, 2020, pp. 1--10.
\end{thebibliography}
\end{document}
